% !TeX root = ./../main.tex

\section{Korekcja, mitygacja i redukcja błędów}\label{sec:tb:qec}
\agentinfo{
Uzgodnienia (PHD-TS-14, 2026-11-26): tutaj definicje i mechanizmy metod; stan badań i wyniki prac w \agentref{chapters/literature_review.tex}.
Terminologia całej pracy: korekcja błędów oznacza wyłącznie QEC, poprawianie wyniku po stronie klasycznej to korekta wyniku, mitygacja bez dopowiedzenia to metody QEM inne niż redukcja.
Metodę PDU definiuje rozdział 4; kryteria K1--K6 i tabela tab:tb03:cost służą dyskusjom w rozdziałach 4--6 i 8.
Notacja: skala szumu c, quasi-prawdopodobieństwa omega, koszt próbkowania makro pecnorm (gamma, kappa, lambda, eta i q są zajęte).
Zaadaptowane źródła własne autora: \agentref{artykuły/magisterka_werner_final/chapters/01.tex} \agentref{artykuły/Quantum_Errors_ICCS/main.tex} \agentref{artykuły/CoDIT_2025_final_paper/root.tex} \agentref{artykuły/Quantum_journal_GAN__Copy_/article.tex}.
}

Błąd sprzętowy, czyli systematyczna składowa błędu wyniku w rozkładzie~\eqref{eq:tb01:errordecomp}, nie maleje ze wzrostem liczby powtórzeń obwodu, więc trzeba mu przeciwdziałać innymi środkami.
W tym podrozdziale zdefiniowano trzy rodziny metod, które temu służą, opisano mechanizmy metod bazowych i wprowadzono kryteria kosztu stosowane w dalszych rozdziałach.
Stan badań i wyniki eksperymentalne omówiono w rozdziale~\ref{chapter:literature_review}.

\subsection{Terminologia przyjęta w pracy}\label{subsec:tb:qec:taxonomy}

W pracy rozróżnia się korekcję, mitygację i redukcję błędów; o przynależności metody decyduje to, czy zmienia ona obwód kwantowy zadania i jakich dodatkowych zasobów wymaga (rys.~\ref{fig:tb03:taxonomy}).

\begin{Definition}[Kwantowa korekcja błędów]\label{def:tb03:qec}
\glsdisp{error-correction}{Kwantową korekcją błędów} (ang. \emph{quantum error correction}, QEC) nazywa się metody, w których stan $k$ \glsdisp{logical-qubit}{kubitów logicznych} jest zakodowany nadmiarowo w stanie $n>k$ kubitów fizycznych, a błędy są wykrywane w trakcie obliczeń pomiarem \glsdisp{syndrome}{syndromu} i odwracane operacją zależną od jego wyniku, bez pomiaru zakodowanej informacji.
\end{Definition}

\begin{Definition}[Mitygacja błędów kwantowych]\label{def:tb03:qem}
\glsdisp{error-mitigation}{Mitygacją błędów kwantowych} (ang. \emph{quantum error mitigation}, QEM) nazywa się metody, które zmniejszają wywołane szumem obciążenie estymatora wartości oczekiwanej lub rozkładu wyników przez klasyczną obróbkę wyników wielu wykonań obwodu, jego zmodyfikowanych wariantów lub obwodów kalibracyjnych, bez kodowania informacji w kubitach logicznych.
\end{Definition}

\begin{Definition}[Redukcja błędów kwantowych]\label{def:tb03:reduction}
\glsdisp{error-reduction}{Redukcją błędów kwantowych} nazywa się w niniejszej pracy metody, które korygują wynik obliczeń kwantowych wyłącznie po stronie klasycznej, modelem korekty wyznaczonym eksperymentalnie dla danego zadania obliczeniowego i urządzenia, przy czym obwód zadania jest wykonywany bez zmian: bez dodatkowych kubitów i bramek, a dodatkowe wykonania służą tylko kalibracji modelu.
\end{Definition}

\begin{figure}[tb]
  \centering
  % !TeX root = ./../../main.tex
% Miejsce ingerencji trzech rodzin metod w proces obliczeń (podrozdział 2.3, rys. fig:tb03:taxonomy).
% Rysunek własny (TikZ): górny wiersz to etapy obliczeń opartych na próbkowaniu,
% trzy pasy poniżej leżą pod etapami, które zmienia korekcja, mitygacja i redukcja błędów.
\begin{tikzpicture}[
  font=\footnotesize,
  stage/.style={draw, rounded corners=2pt, minimum height=0.8cm, text width=3.0cm, align=center, inner sep=2pt},
  band/.style={draw, fill=black!8, inner sep=2.5pt, align=center, minimum height=0.8cm},
  >={Stealth[length=2mm]}]
  % etapy obliczeń
  \node[stage] (circ) at (0,0)    {obwód kwantowy\\zadania};
  \node[stage] (meas) at (3.7,0)  {wykonania obwodu\\i~pomiar};
  \node[stage] (post) at (7.4,0)  {klasyczna obróbka\\statystyk};
  \node[stage] (res)  at (11.1,0) {wynik zadania};
  \draw[->] (circ) -- (meas);
  \draw[->] (meas) -- (post);
  \draw[->] (post) -- (res);
  % pasy: zakres ingerencji metod (szerokość pasa = etapy, które metoda zmienia)
  \node[band, text width=6.6cm] (qec) at (1.85,-1.05)
    {\textbf{korekcja (QEC):} nadmiarowe kubity,\\pomiar syndromu, operacje naprawcze};
  \node[band, text width=10.3cm] (qem) at (3.7,-2.0)
    {\textbf{mitygacja (QEM):} warianty obwodu lub obwody kalibracyjne,\\dodatkowe wykonania, obróbka zebranych statystyk};
  \node[band, text width=6.6cm] (red) at (9.25,-2.95)
    {\textbf{redukcja:} korekta wyniku modelem\\z~kalibracji; obwód zadania bez zmian};
\end{tikzpicture}

  \caption[Miejsce ingerencji korekcji, mitygacji i redukcji błędów w proces obliczeń]{Miejsce ingerencji trzech rodzin metod w proces obliczeń; pas każdej rodziny leży pod etapami, które ta rodzina zmienia.}
  \label{fig:tb03:taxonomy}
\end{figure}

Korekcja usuwa błędy ze stanu kwantowego w każdym wykonaniu obwodu.
Mitygacja i redukcja tego nie robią: w pojedynczym wykonaniu poziom szumu pozostaje taki sam albo wyższy, a poprawia się dopiero wynik wyznaczony ze zbioru wykonań~\cite{cai2023quantum}.
Tytułowe ,,ograniczanie błędów'' obejmuje właśnie te dwie rodziny.
Do redukcji należą metody będące przedmiotem pracy: funkcja odwracania zniekształceń fazy (PDU), w której model korekty wyznacza się z pomiarów funkcji referencyjnej (rozdział~\ref{chapter:pdu}), oraz korekta odtworzonych obrazów modelem generatywnym uczonym na parach obraz zaszumiony--obraz wzorcowy (rozdział~\ref{chapter:dl_error_reduction}).
Termin korekcja błędów jest w pracy zastrzeżony dla QEC, a poprawianie wyniku po stronie klasycznej, będące krokiem metod redukcji, nazywane jest korektą wyniku.

Granice między rodzinami są częściowo umowne.
Mitygacja nie ma powszechnie przyjętej definicji; definicja robocza z przeglądu~\cite{cai2023quantum}, za którą idzie definicja~\ref{def:tb03:qem}, obejmuje każdą obróbkę wyników zbioru wykonań obwodu zmniejszającą obciążenie wywołane szumem.
W tym szerokim sensie redukcja jest szczególnym przypadkiem mitygacji, a angielski termin \emph{error reduction} nie oznacza w literaturze odrębnej klasy metod.
W pracy wyodrębniono ją, ponieważ nie modyfikuje obwodu zadania i jest zorientowana na problem: korektę wyznacza się dla wyniku całego zadania, bez modelu szumu bramek lub pomiaru.
Słowo mitygacja bez dopowiedzenia oznacza dalej metody QEM inne niż redukcja.
Blisko redukcji leży mitygacja błędu odczytu (podrozdział~\ref{subsec:tb:qec:mitigation}), która także nie zmienia obwodu zadania i sprowadza się do kalibracji oraz obliczeń klasycznych~\cite{maciejewski_mitigation_2020,bravyi_mitigating_2021}.
Kalibruje się w niej jednak sam pomiar, niezależnie od obwodu, więc korekta przenosi się na dowolne zadanie, ale obejmuje tylko błąd odczytu; w redukcji kalibruje się wynik całego zadania, więc korekta obejmuje łącznie skutki błędów bramek, dekoherencji i odczytu, lecz nie przenosi się na inne zadania.
Jeszcze bliżej leżą metody, w których korektę wyniku kalibruje się na obwodach o znanym wyniku~\cite{jattana_general_2020,czarnik_error_2021,lolur_reference-state_2023}; ich zestawienie z metodą PDU zawiera podrozdział~\ref{sec:lr:qem} (tabela~\ref{tab:lr:qem:calibrated}).
\agentinfo{
Dopisek orkiestratora (2026-11-27): zdanie o metodach kalibrowanych na obwodach o znanym wyniku dodano po zebraniu literatury do sekcji 3.3, żeby podrozdział nie był sprzeczny z przeglądem literatury.
Decyzja autora z 2026-11-27: PDU jest metodą rozwijaną w ramach tego doktoratu i ma być opisana jako ważny wkład rozprawy, z wyróżnikami podanymi w sekcji 3.3 (korekta indeksowana znanym argumentem i zbadanie stabilności błędu między kalibracjami).
Nadal do decyzji autora: czy redukcję opisywać jako kategorię obejmującą te metody (definicję def:tb03:reduction spełniają także one) oraz czy dodać dla nich wiersz w tabeli tab:tb03:cost, która przypisuje całej mitygacji zorientowanie na bramkę.
}
Poza trzema definicjami pozostają metody pośrednie: weryfikacja symetrii, która stosuje sprawdzenia podobne do pomiaru syndromu, lecz wykonania z wykrytym błędem odrzuca, zamiast je naprawiać, oraz dynamiczne odsprzęganie, technika sterowania działająca bez pomiaru i obróbki wyników, nazywana tłumieniem błędów~\cite{cai2023quantum,ezzell_dynamical_2023}; omówiono je razem z mitygacją.
We wcześniejszych pracach autora metodę PDU nazywano także korekcją błędów~\cite{werner_experiment_driven_2022}, a jej działanie określano jako mitygację~\cite{werner_magisterka_2022}; w rozprawie stosowany jest wyłącznie termin redukcja.

\subsection{Kwantowa korekcja błędów}\label{subsec:tb:qec:codes}

Klasyczne kody korekcyjne powielają bity i rozstrzygają o ich wartości głosowaniem, czego nie da się przenieść wprost na kubity z trzech powodów~\cite{NielsenChuang2010,bib:error_corr_intr}.
Twierdzenie o zakazie klonowania wyklucza istnienie operacji unitarnej $U$ spełniającej $U\ket{\psi}\ket{0}=\ket{\psi}\ket{\psi}$ dla każdego stanu $\ket{\psi}$, więc nieznanego stanu nie można powielić~\cite{wootters_single_1982}.
Pomiar kubitów danych niszczy superpozycję (podrozdział~\ref{subsec:tb:measurement}), a błędy tworzą zbiór ciągły: obok odwrócenia bitu (bramka $\mathrm{X}$) występują odwrócenie fazy ($\mathrm{Z}$), ich złożenie ($\mathrm{Y}$) oraz dowolne kombinacje liniowe tych operatorów.

Kody kwantowe omijają te przeszkody: stan zapisuje się w splątanym stanie wielu kubitów zamiast w kopiach, a mierzy się wyłącznie syndrom, czyli wielkości zależne od błędu, lecz nie od zakodowanych amplitud.
Najprostszym przykładem jest trzykubitowy kod powtórzeniowy, w którym stan $\alpha\ket{0}+\beta\ket{1}$ zapisuje się jako $\alpha\ket{000}+\beta\ket{111}=\alpha\ket{0_L}+\beta\ket{1_L}$, gdzie $\ket{0_L}$ i $\ket{1_L}$ są stanami bazowymi kubitu logicznego.
Syndrom tworzą parzystości sąsiednich par kubitów, czyli wyniki pomiaru operatorów $\mathrm{Z}_1\mathrm{Z}_2$ i $\mathrm{Z}_2\mathrm{Z}_3$, odczytywane za pomocą dwóch kubitów pomocniczych (rys.~\ref{fig:tb03:repetition}).
Odwrócenie jednego bitu przenosi stan do jednej z trzech podprzestrzeni ortogonalnych do podprzestrzeni kodowej, a każdej z nich odpowiada inna wartość syndromu, więc błąd można rozpoznać i odwrócić bramką $\mathrm{X}$ bez poznania $\alpha$ i $\beta$.
Przy niezależnych odwróceniach bitów z prawdopodobieństwem $p$ kod zawodzi dopiero przy co najmniej dwóch błędach, czyli z prawdopodobieństwem $3p^2-2p^3$, mniejszym od $p$ dla $p<1/2$; przed odwróceniem fazy nie chroni.

\begin{figure}[tb]
  \centering
  % !TeX root = ./../../main.tex
% Trzykubitowy kod powtórzeniowy z pomiarem syndromu (podrozdział 2.3, rys. fig:tb03:repetition).
% Rysunek własny (quantikz2): kodowanie, kanał szumu, pomiar parzystości Z1Z2 i Z2Z3
% na dwóch kubitach pomocniczych, operacja naprawcza R sterowana klasycznie wynikiem syndromu.
\begin{quantikz}[row sep={0.55cm,between origins}, column sep=0.34cm, font=\small]
\lstick{$\ket{\psi}$} & \ctrl{1} & \ctrl{2} & \gate[3]{\channel} & \ctrl{3} &          &          &          &          & \gate[3]{R} & \rstick[3]{$\alpha\ket{000}+\beta\ket{111}$} \\
\lstick{$\ket{0}$}    & \targ{}  &          &                    &          & \ctrl{2} & \ctrl{3} &          &          &             & \\
\lstick{$\ket{0}$}    &          & \targ{}  &                    &          &          &          & \ctrl{2} &          &             & \\
\lstick{$\ket{0}$}    &          &          &                    & \targ{}  & \targ{}  &          &          & \meter{} & \ctrl[vertical wire=c]{-1}\setwiretype{c} & \rstick{$s_1$} \\
\lstick{$\ket{0}$}    &          &          &                    &          &          & \targ{}  & \targ{}  & \meter{} & \ctrl[vertical wire=c]{-1}\setwiretype{c} & \rstick{$s_2$}
\end{quantikz}

  \caption[Trzykubitowy kod powtórzeniowy z pomiarem syndromu]{Trzykubitowy kod powtórzeniowy: kodowanie stanu $\ket{\psi}=\alpha\ket{0}+\beta\ket{1}$, kanał szumu $\channel$, pomiar syndromu $(s_1,s_2)$ na kubitach pomocniczych i sterowana nim operacja naprawcza $R$.}
  \label{fig:tb03:repetition}
\end{figure}

Pomiar syndromu rzutuje dowolny błąd pojedynczego kubitu na jeden z operatorów $\identity$, $\mathrm{X}$, $\mathrm{Y}$, $\mathrm{Z}$, dlatego kod korygujący odwrócenia bitu i fazy koryguje każdy błąd jednokubitowy~\cite{NielsenChuang2010}.
Pierwszym takim kodem był dziewięciokubitowy kod Shora~\cite{bib:shor_code}, w którym każdy z trzech kubitów kodu powtórzeniowego chroniącego fazę jest zapisany trzema kubitami kodu chroniącego bit.
Najmniejszy kod o tej własności wykorzystuje pięć kubitów~\cite{bib:perfect_quantum_error_correction_code}: jeżeli brakowi błędu i każdemu z $3n$ błędów jednokubitowych ma odpowiadać osobna para ortogonalnych stanów, to musi zachodzić $2(3n+1)\le2^n$, skąd $n\ge5$.

Wymienione kody należą do rodziny kodów stabilizatorowych~\cite{gottesman_stabilizer_1997,NielsenChuang2010}.
Kod taki określa grupa stabilizatorów $\stabgroup$: przemienna, niezawierająca $-\identity$ podgrupa grupy Pauliego, czyli grupy iloczynów tensorowych operatorów $\identity$, $\mathrm{X}$, $\mathrm{Y}$, $\mathrm{Z}$ na $n$ kubitach z czynnikami $\pm1$, $\pm\mathrm{i}$.
Podprzestrzeń kodową tworzą stany spełniające $S\ket{\psi}=\ket{\psi}$ dla każdego $S\in\stabgroup$; przy $n-k$ niezależnych generatorach grupy mieści ona $k$ kubitów logicznych.
Syndromem jest ciąg wyników pomiaru generatorów: błąd będący operatorem Pauliego zmienia na $-1$ wyniki tych generatorów, z którymi antykomutuje.
Kod oznacza się symbolem $[[n,k,d]]$, w którym dystans $d$ jest najmniejszą liczbą kubitów, na które musi działać operator Pauliego, aby zmienić stan logiczny bez zmiany syndromu; taki kod, na przykład kod Shora $[[9,1,3]]$ albo pięciokubitowy $[[5,1,3]]$, koryguje dowolne błędy na $\lfloor(d-1)/2\rfloor$ kubitach.

W kodach powierzchniowych kubity leżą na dwuwymiarowej siatce, a generatorami są iloczyny operatorów $\mathrm{X}$ albo $\mathrm{Z}$ na co najwyżej czterech sąsiednich kubitach danych~\cite{kitaev_fault-tolerant_2003,bib:surface_codes_fowler}.
Każdy generator mierzy się kubitem pomocniczym sprzężonym tylko z najbliższymi sąsiadami, co sprzyja realizacji na urządzeniach o płaskiej topologii połączeń, a dystans rośnie z rozmiarem siatki.
Pomiary generatorów powtarza się cyklicznie w trakcie obliczeń, a klasyczny dekoder wyznacza z ciągu syndromów najbardziej prawdopodobny układ błędów.

O skalowalności korekcji rozstrzyga twierdzenie progowe: jeżeli prawdopodobieństwo błędu pojedynczej operacji $p$ jest mniejsze od wartości progowej $\pth$, to błąd obliczeń logicznych można uczynić dowolnie małym kosztem narzutu rosnącego polilogarytmicznie z rozmiarem obliczeń~\cite{aharonov_fault-tolerant_2008,knill_resilient_1998,NielsenChuang2010}.
Dla kodu powierzchniowego próg jest rzędu $1\%$ (w modelu błędów z pracy~\cite{bib:surface_codes_fowler} $\pth=0{,}57\%$), a poniżej progu prawdopodobieństwo błędu logicznego maleje wykładniczo z dystansem, $\plog\propto(p/\pth)^{\lfloor(d+1)/2\rfloor}$.
Ceną jest narzut zasobów: potrzebne są kubity pomiarowe, pomiary w trakcie obwodu i dekodowanie w czasie rzeczywistym, a według oszacowania z pracy~\cite{bib:surface_codes_fowler} kubit logiczny o użytecznie małym błędzie wymaga rzędu $10^3$--$10^4$ kubitów fizycznych.
Ograniczeniem jest też topologia: bramki działają tylko między sąsiadami na siatce, więc operacje na odległych kubitach logicznych wymagają dodatkowych kroków i dodatkowego miejsca~\cite{litinski_game_2019}.
Na urządzeniach ery NISQ (podrozdział~\ref{subsec:tb:noise:nisq}) kubity przeznaczone na korekcję pomniejszają i tak niewielką pulę dostępną dla obliczeń, co motywuje metody opisane dalej.

\subsection{Mitygacja błędów}\label{subsec:tb:qec:mitigation}

Metody mitygacji łączy wspólny schemat~\cite{cai2023quantum}: ze statystyk zaszumionych wykonań szacuje się wartość oczekiwaną obserwabli $O$ albo rozkład wyników obwodu idealnego.
Estymator po mitygacji ma mniejsze obciążenie, lecz większą wariancję niż estymator~\eqref{eq:tb01:estimator}, co trzeba skompensować większą liczbą wykonań obwodu.
Ten narzut próbkowania rośnie wykładniczo ze średnią liczbą błędów przypadających na jedno wykonanie obwodu, co ogranicza mitygację do obwodów o umiarkowanym rozmiarze.

\paragraph{Ekstrapolacja do zerowego szumu} (ang. \emph{zero-noise extrapolation}, ZNE)~\cite{richardson-extrapolation,quasi-probability}.
Ten sam obwód wykonuje się przy $m$ poziomach szumu $1=c_1<c_2<\dots<c_m$, gdzie $c=1$ oznacza poziom naturalny, i dla każdego wyznacza wartość $E(c_j)$ obserwabli.
Jeżeli zależność $E(c)$ przybliża wielomian stopnia $m-1$, to jego wartość w zerze, zwana ekstrapolacją Richardsona, wynosi
\begin{equation}\label{eq:tb03:zne}
  \hat{E}(0)=\sum_{j=1}^{m}\beta_j\,E(c_j),\qquad \beta_j=\prod_{l\neq j}\frac{c_l}{c_l-c_j}.
\end{equation}
Szum zwiększa się przez wydłużanie impulsów sterujących~\cite{kandala2021hardware} albo, bez dostępu do impulsów, przez zastępowanie bramki $G$ ciągiem $GG^{\dagger}G$~\cite{giurgica-tiron_digital_2020}.
Metoda zakłada, że szum daje się skalować w kontrolowany sposób, nie zmienia się w czasie zbierania danych i jest na tyle słaby, że $E(c)$ jest funkcją gładką.
Nie wymaga dodatkowych kubitów ani modelu szumu; kosztem jest $m$ serii wykonań, głębsze obwody przy powielaniu bramek oraz wzrost wariancji wyniku (wariancje pomiarów sumują się z wagami $\beta_j^2$).

\paragraph{Probabilistyczne znoszenie błędów} (ang. \emph{probabilistic error cancellation}, PEC), nazywane też metodą quasi-prawdopodobieństw~\cite{quasi-probability,endo_practical_2018}.
Idealną operację $\mathcal{U}$ przedstawia się jako kombinację liniową operacji $\mathcal{O}_i$ wykonalnych na zaszumionym urządzeniu:
\begin{equation}\label{eq:tb03:pec}
  \mathcal{U}=\sum_i\omega_i\,\mathcal{O}_i,\qquad \sum_i\omega_i=1,\qquad \pecnorm=\sum_i|\omega_i|\ge1,
\end{equation}
przy czym część współczynników $\omega_i$ jest ujemna, więc nie tworzą one rozkładu prawdopodobieństwa.
W każdym wykonaniu operację $\mathcal{O}_i$ losuje się z prawdopodobieństwem $|\omega_i|/\pecnorm$, a wynik pomiaru mnoży przez $\pecnorm\operatorname{sgn}(\omega_i)$; średnia tak ważonych wyników jest nieobciążonym estymatorem wartości idealnej.
Rozkład~\eqref{eq:tb03:pec} wymaga znajomości kanału szumu każdej bramki, wyznaczanego tomografią (podrozdział~\ref{subsec:tb:noise:characterisation}), i jego stałości w czasie.
W obwodzie złożonym z wielu bramek współczynniki $\pecnorm$ się mnożą, a liczba wykonań potrzebna do zachowania dokładności rośnie jak kwadrat ich iloczynu, czyli wykładniczo z liczbą bramek.

\paragraph{Mitygacja błędu odczytu macierzą kalibracyjną}~\cite{maciejewski_mitigation_2020,bravyi_mitigating_2021}.
Metoda odwraca klasyczny model błędu odczytu~\eqref{eq:tb:noise:readout}.
Kolumny macierzy $\confmat$ wyznacza się, przygotowując kolejne stany bazowe i mierząc je, a korektę uzyskuje się, działając macierzą $\confmat^{-1}$ na wektor zmierzonych częstości; jeżeli wynik zawiera wartości ujemne, zastępuje się go najbliższym rozkładem prawdopodobieństwa.
Pełna kalibracja rejestru $n$ kubitów wymaga $2^n$ obwodów, a przy założeniu niezależności błędów odczytu poszczególnych kubitów macierz ma tylko $2n$ parametrów.
Metoda zakłada, że błąd przygotowania stanów bazowych jest pomijalny, a macierz nie zmienia się między kalibracją a użyciem.
Obwód zadania pozostaje bez zmian, lecz korygowany jest wyłącznie błąd odczytu.

\paragraph{Dynamiczne odsprzęganie} (ang. \emph{dynamical decoupling}, DD)~\cite{viola_dynamical_1999,ezzell_dynamical_2023}.
W okresach, w których kubit czeka na kolejną bramkę, wstawia się ciąg szybkich impulsów, na przykład $\mathrm{X}$--$\mathrm{X}$ albo $\mathrm{X}$--$\mathrm{Y}$--$\mathrm{X}$--$\mathrm{Y}$, którego złożenie jest identycznością z dokładnością do fazy globalnej, lecz który uśrednia wolnozmienne sprzężenie kubitu z otoczeniem.
Wymaga to szumu zmieniającego się wolno w porównaniu z odstępem między impulsami oraz dokładnych impulsów; metoda dodaje bramki, ale nie wymaga dodatkowych kubitów, wykonań obwodu ani obróbki wyników.

Pozostałe metody mitygacji, w tym weryfikację symetrii i metody uczone na obwodach o znanym wyniku, omówiono w podrozdziale~\ref{sec:lr:qem}.

\subsection{Koszt zasobów jako kryterium porównania}\label{subsec:tb:qec:cost}

Opisane metody płacą za poprawę wyniku zasobami różnego rodzaju, dlatego w rozdziałach~\ref{chapter:pdu}--\ref{chapter:dl_error_reduction} oraz~\ref{chapter:classification} porównuje się je według sześciu kryteriów.
Są to: dodatkowe kubity, nieniosące danych zadania (K1); dodatkowe bramki, zwiększające głębokość obwodu zadania i jego podatność na szum (K2); dodatkowe wykonania obwodu, czyli obwody kalibracyjne i powtórzenia potrzebne do zachowania dokładności statystycznej (K3); koszt obliczeń klasycznych, w tym konieczność prowadzenia ich w czasie rzeczywistym (K4); założenia o szumie, czyli wymagany model szumu i jego stabilność między kalibracją a użyciem (K5; podrozdział~\ref{subsec:tb:noise:sources}); zorientowanie metody (K6).
Metoda zorientowana na bramkę przeciwdziała błędom elementów urządzenia (bramek, kubitów w~spoczynku, pomiaru) i stosuje się do dowolnego obwodu, a metoda zorientowana na problem wyznacza korektę dla wyniku konkretnego zadania, urządzenia i układu kubitów, traktując obwód jak czarną skrzynkę.
Kryteria K1 i K2 określają, jaka część zasobów urządzenia pozostaje dla właściwych obliczeń, co jest istotne w przetwarzaniu obrazów, które wymaga wielu kubitów i głębokich obwodów.
Zestawienie w tabeli~\ref{tab:tb03:cost} pokazuje, że redukcja błędów zajmuje położenie skrajne: poza kalibracją nie zużywa zasobów kwantowych, lecz ma najwęższy zakres stosowalności i opiera się na założeniu powtarzalności błędu, które trzeba sprawdzić eksperymentalnie (rozdział~\ref{chapter:pdu}).

\begin{table}[htb]
  \centering
  \caption[Koszt zasobów metod korekcji, mitygacji i redukcji błędów]{Koszt zasobów metod korekcji, mitygacji i redukcji błędów: dodatkowe kubity (K1), bramki (K2) i wykonania obwodu (K3), obliczenia klasyczne (K4), założenia o szumie (K5), zorientowanie na bramkę albo na problem (K6). Zestawienie na podstawie podrozdziałów~\ref{subsec:tb:qec:codes} i~\ref{subsec:tb:qec:mitigation}.}
  \label{tab:tb03:cost}
  \footnotesize
  \setlength{\tabcolsep}{3pt}
  \newcommand{\tbc}[1]{\raggedright\hspace{0pt}#1}
  \begin{tabular}{@{}p{2.5cm}p{0.6cm}p{2.3cm}p{2.5cm}p{2.3cm}p{2.95cm}p{1.45cm}@{}}
    \toprule
    \tbc{\textbf{Metoda}} & \tbc{\textbf{K1}} & \tbc{\textbf{K2}} & \tbc{\textbf{K3}} & \tbc{\textbf{K4}} & \tbc{\textbf{K5}} & \tbc{\textbf{K6}} \tabularnewline
    \midrule
    \tbc{QEC (kod powierzchniowy)} & \tbc{tak} & \tbc{cykle pomiaru syndromu} & \tbc{nie} & \tbc{dekoder czasu rzeczywistego} & \tbc{błąd operacji poniżej progu} & \tbc{bramka} \tabularnewline
    \tbc{ZNE} & \tbc{nie} & \tbc{przy powielaniu bramek} & \tbc{$m$ serii, wzrost wariancji} & \tbc{dopasowanie krzywej} & \tbc{szum skalowalny, stały, słaby} & \tbc{bramka} \tabularnewline
    \tbc{PEC} & \tbc{nie} & \tbc{wstawiane operacje} & \tbc{rosną wykładniczo} & \tbc{losowanie i~ważenie} & \tbc{znany, stały kanał każdej bramki} & \tbc{bramka} \tabularnewline
    \tbc{Mitygacja błędu odczytu} & \tbc{nie} & \tbc{nie} & \tbc{kalibracja, do $2^n$ obwodów} & \tbc{odwrócenie macierzy} & \tbc{klasyczny, stały błąd odczytu} & \tbc{bramka} \tabularnewline
    \tbc{DD} & \tbc{nie} & \tbc{impulsy w~przerwach} & \tbc{nie} & \tbc{brak} & \tbc{szum wolnozmienny} & \tbc{bramka} \tabularnewline
    \tbc{Redukcja (rozdz.~\ref{chapter:pdu},~\ref{chapter:dl_error_reduction})} & \tbc{nie} & \tbc{nie} & \tbc{kalibracja na zadaniach wzorcowych} & \tbc{model korekty} & \tbc{powtarzalny błąd zadania} & \tbc{problem} \tabularnewline
    \bottomrule
  \end{tabular}
\end{table}
